\chapter{Создание макета интернет-магазина}
\label{ch:practice}

\section{Задание~1. Главная страница (десктоп~1440\,px)}

\subsection{Шапка с меню и иконками}

В верхней части фрейма \texttt{Desktop~1440} создаётся горизонтальный
блок шапки высотой ~80\,px.
Состав:
\begin{itemize}
  \item \textbf{Логотип} (слева) --- текстовый слой с названием магазина
        (например, \textit{Greenly}) или иконка листа из
        \texttt{Figma~Community};
  \item \textbf{Навигация} (по центру или справа) --- горизонтальный
        список текстовых слоёв: <<Каталог>>, <<Доставка>>, <<Уход>>,
        <<О~нас>>, <<Контакты>>;
  \item \textbf{Иконки} (справа) --- поиск, избранное, корзина, профиль.
        Иконки берутся из \texttt{Figma~Community} (плагины
        \texttt{Iconify}, \texttt{Material~Icons}) или ищутся в
        интернете в формате \texttt{SVG}.
\end{itemize}
Все элементы шапки объединяются во фрейм \texttt{Header} (\texttt{Ctrl+Alt+G}),
что позволит позже применить к нему привязки.

% — Скриншот: шапка с меню и иконками —
\begin{figure}[h]
  \centering
  \includegraphics[width=0.85\textwidth]{img/03_header}
  \caption{Шапка интернет-магазина с навигацией и иконками}
  \label{fig:header}
\end{figure}

\subsection{Обложка (Hero-блок)}

Hero-блок --- крупный заголовочный экран сразу под шапкой. Состав:
\begin{itemize}
  \item фоновое изображение растения (\texttt{Unsplash}, ключевые слова
        \textit{houseplant}, \textit{indoor~plants}) либо
        одноцветная/градиентная заливка;
  \item крупный заголовок (\texttt{H1}, \texttt{60--72\,pt}, жирное
        начертание) --- слоган магазина;
  \item подзаголовок (\texttt{18--20\,pt}) --- краткое описание;
  \item CTA-кнопка \textit{<<В каталог>>}: прямоугольник с радиусом
        \texttt{8\,px}, заливка контрастного цвета, текстовый слой
        внутри.
\end{itemize}

% — Скриншот: Hero-блок с заголовком и CTA-кнопкой —
\begin{figure}[h]
  \centering
  \includegraphics[width=0.85\textwidth]{img/04_hero}
  \caption{Обложка с заголовком, описанием и CTA-кнопкой}
  \label{fig:hero}
\end{figure}

\subsection{Блоки с фотографиями}

Под Hero-блоком располагается секция с маленькими тематическими
карточками (3--4 в ряд). Каждая карточка --- фрейм со скруглением углов
\(12\,\)px, содержащий:
\begin{itemize}
  \item изображение из \texttt{Unsplash} (заливка прямоугольника);
  \item краткую подпись над/под изображением.
\end{itemize}
Карточки можно использовать как навигационные превью категорий
(<<Суккуленты>>, <<Папоротники>>, <<Цветущие>>, <<Тропические>>).

% — Скриншот: блок с фотографиями растений —
\begin{figure}[h]
  \centering
  \includegraphics[width=0.85\textwidth]{img/05_gallery}
  \caption{Блок с тематическими карточками-категориями}
  \label{fig:gallery}
\end{figure}

\textbf{Самопроверка задания~1:}
\begin{itemize}
  \item меню в шапке создано;
  \item собрана обложка с текстом и кнопкой;
  \item добавлены блоки с маленькими фото.
\end{itemize}

\section{Задание~2. Дополнительные блоки и адаптивные версии}

\subsection{Блок <<Предложения месяца>>}

Блок представляет собой каталог из 2~рядов по 3--4 карточки товара.

Структура одной карточки (фрейм \texttt{Card}, ширина ~280\,px,
высота ~360\,px):
\begin{enumerate}
  \item изображение товара сверху (\(280\times240\,\)px);
  \item название товара (\texttt{16\,pt}, жирное);
  \item краткое описание / латинское название (\texttt{12\,pt},
        тонкое);
  \item цена (\texttt{18\,pt}) и кнопка \textit{<<В корзину>>}.
\end{enumerate}
Все элементы карточки объединяются во фрейм
(\texttt{Ctrl+Alt+G}). Привязки внутри карточки настраиваются так,
чтобы при изменении её ширины фото и текст вели себя предсказуемо
(фото --- \texttt{Left~and~Right} + \texttt{Top}, цена --- \texttt{Left}
+ \texttt{Bottom}).

После завершения первой карточки --- копирование её
(\texttt{Ctrl+D}) и замена текста/изображения для остальных. Второй
ряд получается копированием первого ряда целиком.

% — Скриншот: блок «Предложения месяца» с карточками товаров —
\begin{figure}[h]
  \centering
  \includegraphics[width=0.85\textwidth]{img/06_offers}
  \caption{Блок <<Предложения месяца>> --- два ряда карточек товара}
  \label{fig:offers}
\end{figure}

\subsection{Футер}

Футер --- нижний блок страницы. Состав:
\begin{itemize}
  \item заголовок \textit{<<Подпишись на новости>>};
  \item поле ввода e-mail (прямоугольник с обводкой и плейсхолдер-текстом);
  \item кнопка \textit{<<Подписаться>>};
  \item ниже --- три колонки ссылок: <<Магазин>>, <<Информация>>,
        <<Помощь>>;
  \item иконки социальных сетей (Telegram, ВКонтакте, YouTube,
        Instagram) --- 4--5 шт., из \texttt{Figma~Community};
  \item копирайт-строка \textit{<<\copyright~\cfgYear~Greenly>>}.
\end{itemize}

% — Скриншот: футер с формой подписки и иконками соцсетей —
\begin{figure}[h]
  \centering
  \includegraphics[width=0.85\textwidth]{img/07_footer}
  \caption{Футер с подпиской, ссылками и иконками социальных сетей}
  \label{fig:footer}
\end{figure}

\subsection{Layout grid}

Колоночная сетка применяется к корневому фрейму \texttt{Desktop~1440}:
\begin{itemize}
  \item тип \texttt{Columns};
  \item количество --- \texttt{12};
  \item \texttt{Width} --- \texttt{Auto} (\texttt{Stretch});
  \item \texttt{Margin} --- \texttt{80\,px};
  \item \texttt{Gutter} --- \texttt{24\,px};
  \item видимость сетки переключается \texttt{Ctrl+G}.
\end{itemize}
После применения сетки все ключевые блоки (шапка, Hero, карточки,
футер) выравниваются по колонкам.

% — Скриншот: десктоп-макет с включённой 12-колоночной сеткой —
\begin{figure}[h]
  \centering
  \includegraphics[width=0.85\textwidth]{img/08_grid}
  \caption{12-колоночная \texttt{Layout grid} на фрейме \texttt{Desktop~1440}}
  \label{fig:grid}
\end{figure}

\subsection{Constraints}

Для каждого ключевого блока выбираются привязки:
\begin{itemize}
  \item \texttt{Header}, \texttt{Footer} --- \texttt{Left~and~Right}
        + соответственно \texttt{Top}/\texttt{Bottom};
  \item Hero-фон --- \texttt{Left~and~Right} + \texttt{Top};
  \item ряд карточек --- \texttt{Left~and~Right} + \texttt{Top}; каждая
        карточка внутри --- \texttt{Scale} по горизонтали;
  \item иконки в шапке --- \texttt{Right} + \texttt{Top};
  \item логотип --- \texttt{Left} + \texttt{Top}.
\end{itemize}

% — Скриншот: панель Constraints для карточки товара —
\begin{figure}[h]
  \centering
  \includegraphics[width=0.85\textwidth]{img/09_constraints}
  \caption{Панель \texttt{Constraints} для карточки товара}
  \label{fig:constraints}
\end{figure}

\subsection{Адаптивная версия~768\,px (планшет)}

Алгоритм:
\begin{enumerate}
  \item Скопировать корневой фрейм (\texttt{Ctrl+D}),
        переименовать в \texttt{Tablet~768}.
  \item На правой панели в разделе \texttt{Design} изменить ширину
        фрейма с \(1440\) на \(768\,\)px. Благодаря привязкам Figma
        автоматически адаптирует контент.
  \item Поправить детали вручную:
  \begin{itemize}
    \item перестроить карточки товара в 2~колонки вместо 3--4;
    \item уменьшить отступы шапки и Hero;
    \item обновить \texttt{Layout grid}: \texttt{6}~колонок,
          \texttt{Margin~32\,px}, \texttt{Gutter~16\,px}.
  \end{itemize}
\end{enumerate}

% — Скриншот: адаптивная версия 768 px —
\begin{figure}[h]
  \centering
  \includegraphics[width=0.85\textwidth]{img/10_tablet_768}
  \caption{Адаптивная версия макета шириной \(768\,\)px (планшет)}
  \label{fig:tablet}
\end{figure}

\subsection{Адаптивная версия~360\,px (мобильный)}

Алгоритм:
\begin{enumerate}
  \item Скопировать \texttt{Tablet~768}, переименовать в
        \texttt{Mobile~360}.
  \item Изменить ширину до \(360\,\)px.
  \item Обработать сужения:
  \begin{itemize}
    \item скрыть горизонтальное меню, заменив его на иконку
          <<гамбургер>>;
    \item уменьшить кегль заголовков (\texttt{H1: 36--40\,pt} вместо
          \texttt{60--72\,pt}; \texttt{H2: 24\,pt});
    \item построить карточки в одну колонку;
    \item перейти на 4-колоночную \texttt{Layout grid}: \texttt{Margin
          16\,px}, \texttt{Gutter 8\,px}.
  \end{itemize}
\end{enumerate}

% — Скриншот: адаптивная версия 360 px —
\begin{figure}[h]
  \centering
  \includegraphics[width=0.85\textwidth]{img/11_mobile_360}
  \caption{Адаптивная версия макета шириной \(360\,\)px (мобильный)}
  \label{fig:mobile}
\end{figure}

\section{Готовый макет в Figma}

Ссылка на Figma-файл с тремя адаптивными версиями макета:

\begin{lstlisting}
TODO: вставить URL общего доступа к Figma-файлу
\end{lstlisting}

\textbf{Самопроверка задания~2:}
\begin{itemize}
  \item создан блок <<Предложения месяца>> и футер;
  \item к корневому фрейму применена 12-колоночная \texttt{Layout grid};
  \item для всех ключевых элементов настроены \texttt{Constraints};
  \item созданы две адаптивные версии (\(768\,\)px и \(360\,\)px).
\end{itemize}

\endinput
